% --- Archon physics formalization source begin ---
% archon:physics
% archon:covers IPhO2026Problems/problem_IPhO_2026_1_A_1.lean
% archon:source-report reports/ipho_2026_k3/problem_IPhO_2026_1_A_1.source.json
% archon:problem-id IPhO_2026_1
% archon:part-id A.1

\chapter{Physics problem IPhO\_2026\_1\_A\_1}
\label{ch:IPhO2026Problems_problem_IPhO_2026_1_A_1}

\paragraph{Problem source.}
Two water reservoirs are separated by a vertical wall MN.  A square slot of
vertical size a*sqrt(2)/2 is sealed by a fully submerged solid cube of side a
and density 3*rho\_0, where rho\_0 is the density of water.  The cube is hinged
frictionlessly at O and may rotate about an axis perpendicular to the figure.
The maximum permitted difference in water levels is Delta h = 1.41 m.  Use
Figure 1a on the source page for the exact geometry and lever arms.

Current subquestion:
Calculate the side length a that makes Delta h = 1.41 m the maximum permissible water-level difference.

\paragraph{Current subquestion.}
Calculate the side length a that makes Delta h = 1.41 m the maximum permissible water-level difference.

\paragraph{Recorded answer/context.}
a = Delta h/(2*sqrt(2)) = 0.50 m.

\paragraph{Figure/image path.}
/root/proposal\_for\_physic/science-mango/ipho\_2026\_source/image/T1\_page-1.png

\paragraph{Formalization target.}
create a compiling Lean file with sorry bodies at `IPhO2026Problems/problem\_IPhO\_2026\_1\_A\_1.lean`.
The Lean declarations must preserve the physical quantities, dimensions or dimensional roles, figure labels, governing-law hypotheses, and final relation expressed by this problem.
Use Mathlib/Physlib names found through LeanExplore where available. If a domain API is missing, introduce faithful local abstractions rather than scalar placeholder aliases.

\begin{theorem}[Physics formalization target]
\label{thm:physics:IPhO_2026_1_A_1:target}
\uses{thm:IPhO2026Problems_problem_IPhO_2026_1_A_1:hydrostatic_gate_side_length_a_target}
The assigned autoformalize agent should translate this physics problem into Lean declarations in the covered file, with theorem and lemma proof bodies written as `by sorry`.
\end{theorem}
\begin{proof}
This is an autoformalization task, not a proof task. Produce faithful statements that can later be proved without weakening the source contract.
\end{proof}
% NOTE: PhysLean-coverage exemption (planner-recorded, iter-002): PhysLean has no hydrostatics/fluid-statics module; nearest hits are Path.topological lemmas (see task\_results/physics-grounding-IPhO2026Problems\_problem\_IPhO\_2026\_1\_A\_1.md). Self-containment is kept with the `import Mathlib` baseline; no irrelevant Physlib import is added. This documents the resolved import policy (iter-001 review ruling) for the blueprint-doctor `missing-physlib-import` check.
% --- Archon physics formalization source end ---

% --- Archon declaration ledger (added iter-008, helper-layer coverage) ---
% One block per live declaration of the covered file, restated first-hand
% from the Lean source; \uses{} mirrors the real logical dependencies.

\section*{Declaration ledger (autoformalization contracts)}

\subsection*{Setup: parameters and figure constants}

\begin{definition}[Transverse figure plane]
\label{def:IPhO2026Problems_problem_IPhO_2026_1_A_1:GatePlane}
\lean{IPhO2026_1_A_1.GatePlane}
The transverse plane of Figure 1a: the gate apparatus is translationally
invariant along the hinge axis, so the configuration lives in a Euclidean
2-space over $\mathbb{R}$ and lever arms are intrinsic vectors of this plane.
\end{definition}
\begin{proof}
Definition; no claim.
\end{proof}

\begin{definition}[Water density]
\label{def:IPhO2026Problems_problem_IPhO_2026_1_A_1:rho0}
\lean{IPhO2026_1_A_1.rho0}
The density $\rho_0$ of water ($\mathrm{kg/m^3}$), carried as an abstract
scalar parameter of the model.
\end{definition}
\begin{proof}
Definition; no claim.
\end{proof}

\begin{definition}[Cube side length]
\label{def:IPhO2026Problems_problem_IPhO_2026_1_A_1:a}
\lean{IPhO2026_1_A_1.a}
The side length $a$ of the cubic gate block (m). Its critical value is the
object of this subquestion, so $a$ itself stays an abstract datum and no
hypothesis fixes it.
\end{definition}
\begin{proof}
Definition; no claim.
\end{proof}

\begin{definition}[Level-difference bound]
\label{def:IPhO2026Problems_problem_IPhO_2026_1_A_1:DeltaH}
\lean{IPhO2026_1_A_1.DeltaH}
The maximum permissible difference $\Delta h$ of the two free-surface levels
(m); the design fixes $\Delta h = 1.41$.
\end{definition}
\begin{proof}
Definition; no claim.
\end{proof}

\begin{definition}[Gravitational field magnitude]
\label{def:IPhO2026Problems_problem_IPhO_2026_1_A_1:g}
\lean{IPhO2026_1_A_1.g}
The magnitude $g$ of the uniform downward gravitational field
($\mathrm{m/s^2}$); it cancels in the final balance, but the weight law and
the hydrostatic law are stated through it.
\end{definition}
\begin{proof}
Definition; no claim.
\end{proof}

\begin{definition}[Positivity regime]
\label{def:IPhO2026Problems_problem_IPhO_2026_1_A_1:PhysicalParameters}
\lean{IPhO2026_1_A_1.PhysicalParameters}
\uses{def:IPhO2026Problems_problem_IPhO_2026_1_A_1:rho0, def:IPhO2026Problems_problem_IPhO_2026_1_A_1:a, def:IPhO2026Problems_problem_IPhO_2026_1_A_1:DeltaH, def:IPhO2026Problems_problem_IPhO_2026_1_A_1:g}
The physical regime of the problem: $0 < \rho_0$, $0 < a$, $0 < \Delta h$ and
$0 < g$. Every later hypothesis bundle extends this regime.
\end{definition}
\begin{proof}
Definition; no claim.
\end{proof}

\begin{definition}[Cube mass]
\label{def:IPhO2026Problems_problem_IPhO_2026_1_A_1:cubeMass}
\lean{IPhO2026_1_A_1.cubeMass}
\uses{def:IPhO2026Problems_problem_IPhO_2026_1_A_1:rho0, def:IPhO2026Problems_problem_IPhO_2026_1_A_1:a}
The mass $m_c = 3\rho_0 a^3$ of the cubic block (kg): volume $a^3$ times the
material density $3\rho_0$ (the cube is three times as dense as water).
\end{definition}
\begin{proof}
Definition; no claim.
\end{proof}

\begin{definition}[Displaced water mass]
\label{def:IPhO2026Problems_problem_IPhO_2026_1_A_1:displacedWaterMass}
\lean{IPhO2026_1_A_1.displacedWaterMass}
\uses{def:IPhO2026Problems_problem_IPhO_2026_1_A_1:rho0, def:IPhO2026Problems_problem_IPhO_2026_1_A_1:a}
The mass $\rho_0 a^3$ of water displaced by the fully submerged cube (kg).
\end{definition}
\begin{proof}
Definition; no claim.
\end{proof}

\begin{definition}[Slot vertical size]
\label{def:IPhO2026Problems_problem_IPhO_2026_1_A_1:slotVerticalSize}
\lean{IPhO2026_1_A_1.slotVerticalSize}
\uses{def:IPhO2026Problems_problem_IPhO_2026_1_A_1:a}
Figure-1a readout (the $45^\circ$ diamond seal): with the cube's sides at
$45^\circ$ to the horizontal, the square slot in the vertical wall MN has
vertical size $a\sqrt{2}/2$, the horizontal thickness of the cube.
\end{definition}
\begin{proof}
Definition; no claim.
\end{proof}

\begin{definition}[Frictionless hinge axis]
\label{def:IPhO2026Problems_problem_IPhO_2026_1_A_1:HingeAxis}
\lean{IPhO2026_1_A_1.HingeAxis}
\uses{def:IPhO2026Problems_problem_IPhO_2026_1_A_1:GatePlane}
Figure content, assumed: the frictionless hinge at the cube's bottom vertex O,
recorded as a fixed point of the figure plane; the plane of the figure passes
through O perpendicular to the hinge axis.
\end{definition}
\begin{proof}
Definition; no claim.
\end{proof}

\subsection*{Governing laws (assumption side)}

\begin{definition}[Weight law]
\label{def:IPhO2026Problems_problem_IPhO_2026_1_A_1:IsWeightForce}
\lean{IPhO2026_1_A_1.IsWeightForce}
\uses{def:IPhO2026Problems_problem_IPhO_2026_1_A_1:g}
Governing law, assumed: a body of mass $m$ in the uniform field experiences
the downward weight force $F = m \cdot g$.
\end{definition}
\begin{proof}
Definition; no claim.
\end{proof}

\begin{definition}[Uniform gravity field]
\label{def:IPhO2026Problems_problem_IPhO_2026_1_A_1:IsUniformGravityField}
\lean{IPhO2026_1_A_1.IsUniformGravityField}
\uses{def:IPhO2026Problems_problem_IPhO_2026_1_A_1:IsWeightForce}
Governing law, assumed: the field strength is the same scalar $g$ at every
point of the apparatus, so two bodies of equal mass receive the same weight
wherever they are placed.
\end{definition}
\begin{proof}
Definition; no claim.
\end{proof}

\begin{definition}[Hydrostatic pressure law]
\label{def:IPhO2026Problems_problem_IPhO_2026_1_A_1:IsHydrostaticPressure}
\lean{IPhO2026_1_A_1.IsHydrostaticPressure}
\uses{def:IPhO2026Problems_problem_IPhO_2026_1_A_1:g}
Governing law, assumed: at depth $d$ below a free surface in a liquid of
density $\rho$ the gauge pressure obeys the linear depth law
$p = \rho \cdot g \cdot d$.
\end{definition}
\begin{proof}
Definition; no claim.
\end{proof}

\begin{definition}[Archimedes buoyancy law]
\label{def:IPhO2026Problems_problem_IPhO_2026_1_A_1:IsBuoyantForce}
\lean{IPhO2026_1_A_1.IsBuoyantForce}
\uses{def:IPhO2026Problems_problem_IPhO_2026_1_A_1:g}
Governing law, assumed: a fully submerged body experiences the upward buoyant
force $B = \rho_0 \cdot V \cdot g$, the weight of the displaced water.
\end{definition}
\begin{proof}
Definition; no claim.
\end{proof}

\begin{definition}[Net immersed weight law]
\label{def:IPhO2026Problems_problem_IPhO_2026_1_A_1:IsNetImmersedWeight}
\lean{IPhO2026_1_A_1.IsNetImmersedWeight}
\uses{def:IPhO2026Problems_problem_IPhO_2026_1_A_1:IsWeightForce, def:IPhO2026Problems_problem_IPhO_2026_1_A_1:IsBuoyantForce, def:IPhO2026Problems_problem_IPhO_2026_1_A_1:cubeMass, def:IPhO2026Problems_problem_IPhO_2026_1_A_1:rho0, def:IPhO2026Problems_problem_IPhO_2026_1_A_1:a}
Governing law, assumed: the effective vertical force carried by the block is
$F = W - B$, where $W = m_c g$ is the cube's weight and
$B = \rho_0 a^3 g$ is the buoyancy on the displaced volume $a^3$.
\end{definition}
\begin{proof}
Definition; no claim.
\end{proof}

\begin{definition}[Weight lever arm at $45^\circ$]
\label{def:IPhO2026Problems_problem_IPhO_2026_1_A_1:weightHorizontalLeverArm}
\lean{IPhO2026_1_A_1.weightHorizontalLeverArm}
\uses{def:IPhO2026Problems_problem_IPhO_2026_1_A_1:a}
Figure-1a readout: the cube's centre lies $a/2$ above O along the vertical
diagonal and the sides lie at $\pi/4$, so the horizontal lever arm of the
immersed weight about O is $(a/2)\sin(\pi/4)$.
\end{definition}
\begin{proof}
Definition; no claim.
\end{proof}

\begin{definition}[Pressure moment readout]
\label{def:IPhO2026Problems_problem_IPhO_2026_1_A_1:PressureMomentReadout}
\lean{IPhO2026_1_A_1.PressureMomentReadout}
\uses{def:IPhO2026Problems_problem_IPhO_2026_1_A_1:DeltaH, def:IPhO2026Problems_problem_IPhO_2026_1_A_1:slotVerticalSize, def:IPhO2026Problems_problem_IPhO_2026_1_A_1:rho0, def:IPhO2026Problems_problem_IPhO_2026_1_A_1:g, def:IPhO2026Problems_problem_IPhO_2026_1_A_1:a}
Figure-1a readout combined with the hydrostatic law, assumed: in the critical
wetted regime both faces of the cube stay fully wetted,
$\Delta h \le a\sqrt{2}/2$, and the net driving couple about O per unit
hinge-axis length is
$\tau = \rho_0 g\, \Delta h \cdot (a\sqrt{2}/2) \cdot (a/2) \cdot
(a\sqrt{2}/4)$, the head times the slot's vertical size times the figure
lever arms, with $a\sqrt{2}/4$ the horizontal arm read off the figure.
\end{definition}
\begin{proof}
Definition; no claim.
\end{proof}

\begin{definition}[Critical torque balance about O]
\label{def:IPhO2026Problems_problem_IPhO_2026_1_A_1:IsCriticalTorqueBalance}
\lean{IPhO2026_1_A_1.IsCriticalTorqueBalance}
\uses{def:IPhO2026Problems_problem_IPhO_2026_1_A_1:weightHorizontalLeverArm, def:IPhO2026Problems_problem_IPhO_2026_1_A_1:slotVerticalSize, def:IPhO2026Problems_problem_IPhO_2026_1_A_1:rho0, def:IPhO2026Problems_problem_IPhO_2026_1_A_1:g, def:IPhO2026Problems_problem_IPhO_2026_1_A_1:DeltaH, def:IPhO2026Problems_problem_IPhO_2026_1_A_1:a}
Governing equilibrium law, assumed: the cube is static about the frictionless
hinge, so the sum of moments about O vanishes — the anticlockwise restoring
moment $F \cdot (a/2)\sin(\pi/4)$ of the net immersed weight balances the
clockwise hydrostatic couple
$\rho_0 g\, \Delta h \cdot (a\sqrt{2}/2) \cdot (a/2) \cdot (a\sqrt{2}/4)$ at
the maximum permissible $\Delta h$.
\end{definition}
\begin{proof}
Definition; no claim.
\end{proof}

\begin{definition}[Full gate assumption bundle]
\label{def:IPhO2026Problems_problem_IPhO_2026_1_A_1:HydrostaticGateSetup}
\lean{IPhO2026_1_A_1.HydrostaticGateSetup}
\uses{def:IPhO2026Problems_problem_IPhO_2026_1_A_1:PhysicalParameters, def:IPhO2026Problems_problem_IPhO_2026_1_A_1:IsUniformGravityField, def:IPhO2026Problems_problem_IPhO_2026_1_A_1:IsHydrostaticPressure, def:IPhO2026Problems_problem_IPhO_2026_1_A_1:PressureMomentReadout, def:IPhO2026Problems_problem_IPhO_2026_1_A_1:IsCriticalTorqueBalance, def:IPhO2026Problems_problem_IPhO_2026_1_A_1:IsNetImmersedWeight}
The complete assumption side of the problem: the positive parameter regime,
the uniform gravity field, the hydrostatic law holding at the hinge O for
some head $h_O$, the Figure-1a pressure readout on the wetted faces, and the
existence of a critically torque-balanced configuration carrying the net
immersed weight. The official answer closed form appears nowhere among these
fields.
\end{definition}
\begin{proof}
Definition; no claim.
\end{proof}

\subsection*{Derivation chain to the official value}

\begin{definition}[Net immersed weight magnitude]
\label{def:IPhO2026Problems_problem_IPhO_2026_1_A_1:netImmersedWeight}
\lean{IPhO2026_1_A_1.netImmersedWeight}
\uses{def:IPhO2026Problems_problem_IPhO_2026_1_A_1:cubeMass, def:IPhO2026Problems_problem_IPhO_2026_1_A_1:displacedWaterMass, def:IPhO2026Problems_problem_IPhO_2026_1_A_1:g}
The derived force magnitude $F = (m_c - \rho_0 a^3)\, g$ carried on the
restoring side of the moment balance: weight minus buoyancy.
\end{definition}
\begin{proof}
Definition; no claim.
\end{proof}

\begin{definition}[Restoring moment about O]
\label{def:IPhO2026Problems_problem_IPhO_2026_1_A_1:restoringMoment}
\lean{IPhO2026_1_A_1.restoringMoment}
\uses{def:IPhO2026Problems_problem_IPhO_2026_1_A_1:netImmersedWeight, def:IPhO2026Problems_problem_IPhO_2026_1_A_1:weightHorizontalLeverArm}
The anticlockwise restoring moment of the net immersed weight about O: the
force times its horizontal lever arm.
\end{definition}
\begin{proof}
Definition; no claim.
\end{proof}

\begin{definition}[Pressure couple magnitude]
\label{def:IPhO2026Problems_problem_IPhO_2026_1_A_1:pressureCoupleMagnitude}
\lean{IPhO2026_1_A_1.pressureCoupleMagnitude}
\uses{def:IPhO2026Problems_problem_IPhO_2026_1_A_1:rho0, def:IPhO2026Problems_problem_IPhO_2026_1_A_1:g, def:IPhO2026Problems_problem_IPhO_2026_1_A_1:DeltaH, def:IPhO2026Problems_problem_IPhO_2026_1_A_1:slotVerticalSize, def:IPhO2026Problems_problem_IPhO_2026_1_A_1:a}
The magnitude of the hydrostatic driving couple about O per unit axis length
in the critical regime,
$\rho_0 g\, \Delta h \cdot (a\sqrt{2}/2) \cdot (a/2) \cdot (a\sqrt{2}/4)$ —
the Figure-1a readout side of the balance law (the slot factor $a\sqrt{2}/2$
and the arm $a\sqrt{2}/4$ are figure readouts).
\end{definition}
\begin{proof}
Definition; no claim.
\end{proof}

\begin{lemma}[Step 1: mass--displacement law]
\label{lem:IPhO2026Problems_problem_IPhO_2026_1_A_1:net_immersed_weight_eq}
\lean{IPhO2026_1_A_1.net_immersed_weight_eq}
\uses{def:IPhO2026Problems_problem_IPhO_2026_1_A_1:netImmersedWeight, def:IPhO2026Problems_problem_IPhO_2026_1_A_1:cubeMass, def:IPhO2026Problems_problem_IPhO_2026_1_A_1:displacedWaterMass}
Under positivity of the parameters the net immersed weight equals
$2\rho_0 a^3 g$.
\end{lemma}
\begin{proof}
Substituting $m_c = 3\rho_0 a^3$ and the displaced mass $\rho_0 a^3$, the
difference $3\rho_0 a^3 - \rho_0 a^3$ collects to $2\rho_0 a^3$; multiplying
by $g$ gives the claim.
\end{proof}

\begin{lemma}[Step 2: lever arm at $45^\circ$]
\label{lem:IPhO2026Problems_problem_IPhO_2026_1_A_1:weight_lever_arm_eq}
\lean{IPhO2026_1_A_1.weight_lever_arm_eq}
\uses{def:IPhO2026Problems_problem_IPhO_2026_1_A_1:weightHorizontalLeverArm}
The horizontal lever arm of the immersed weight about O equals
$a/(2\sqrt{2})$.
\end{lemma}
\begin{proof}
The sine special value $\sin(\pi/4) = \sqrt{2}/2$ gives
$(a/2)(\sqrt{2}/2)$, which equals $a/(2\sqrt{2})$ since
$\sqrt{2}\cdot\sqrt{2} = 2$.
\end{proof}

\begin{lemma}[Step 3: restoring moment]
\label{lem:IPhO2026Problems_problem_IPhO_2026_1_A_1:restoring_moment_eq}
\lean{IPhO2026_1_A_1.restoring_moment_eq}
\uses{def:IPhO2026Problems_problem_IPhO_2026_1_A_1:restoringMoment, lem:IPhO2026Problems_problem_IPhO_2026_1_A_1:net_immersed_weight_eq, lem:IPhO2026Problems_problem_IPhO_2026_1_A_1:weight_lever_arm_eq}
The restoring moment about O equals $\rho_0 g a^4/\sqrt{2}$.
\end{lemma}
\begin{proof}
Multiply the results of Steps 1 and 2,
$2\rho_0 a^3 g \cdot a/(2\sqrt{2})$, and cancel the factor $2$.
\end{proof}

\begin{lemma}[Step 4: hydrostatic couple]
\label{lem:IPhO2026Problems_problem_IPhO_2026_1_A_1:pressure_couple_eq}
\lean{IPhO2026_1_A_1.pressure_couple_eq}
\uses{def:IPhO2026Problems_problem_IPhO_2026_1_A_1:pressureCoupleMagnitude, def:IPhO2026Problems_problem_IPhO_2026_1_A_1:slotVerticalSize}
The driving pressure couple about O simplifies to
$\rho_0 g\, \Delta h\, a^3/4$.
\end{lemma}
\begin{proof}
Insert the slot's vertical size $a\sqrt{2}/2$ into the readout expression and
collect $\sqrt{2}\cdot\sqrt{2} = 2$: the product
$\rho_0 g\, \Delta h \cdot (a\sqrt{2}/2) \cdot (a/2) \cdot (a\sqrt{2}/4)$
becomes $\rho_0 g\, \Delta h\, a^3 \cdot 2/16 = \rho_0 g\, \Delta h\, a^3/4$.
\end{proof}

\begin{lemma}[Step 5: critical balance equation]
\label{lem:IPhO2026Problems_problem_IPhO_2026_1_A_1:critical_balance_eq}
\lean{IPhO2026_1_A_1.critical_balance_eq}
\uses{lem:IPhO2026Problems_problem_IPhO_2026_1_A_1:restoring_moment_eq, lem:IPhO2026Problems_problem_IPhO_2026_1_A_1:pressure_couple_eq}
In the critical configuration, where the restoring moment equals the pressure
couple, the balance reduces to the scalar equation
$\rho_0 g a^4/\sqrt{2} = \rho_0 g\, \Delta h\, a^3/4$.
\end{lemma}
\begin{proof}
Rewrite the left and right sides of the moment-balance hypothesis by the
closed evaluations of Steps 3 and 4.
\end{proof}

\begin{lemma}[Step 6: couple-position trace]
\label{lem:IPhO2026Problems_problem_IPhO_2026_1_A_1:pressure_couple_position_trace}
\lean{IPhO2026_1_A_1.pressure_couple_position_trace}
\uses{def:IPhO2026Problems_problem_IPhO_2026_1_A_1:slotVerticalSize}
The arm of the pressure couple is one quarter of the figure readout
$a\sqrt{2}$: $4 \cdot (a\sqrt{2}/4) = a\sqrt{2}$. A purely geometrical
Figure-1a trace of where the couple acts, consistently with the top face
sitting flush at the slot's upper lip in the critical configuration.
\end{lemma}
\begin{proof}
Four quarters make a whole: $4 \cdot (x/4) = x$ for $x = a\sqrt{2}$ by direct
rational arithmetic.
\end{proof}

\begin{lemma}[Step 7: solving for the side length]
\label{lem:IPhO2026Problems_problem_IPhO_2026_1_A_1:side_length_eq_delta_h_over}
\lean{IPhO2026_1_A_1.side_length_eq_delta_h_over}
\uses{lem:IPhO2026Problems_problem_IPhO_2026_1_A_1:critical_balance_eq}
From the critical balance equation
$\rho_0 g a^4/\sqrt{2} = \rho_0 g\, \Delta h\, a^3/4$ and the positivity of
the parameters, $a = \Delta h/(2\sqrt{2})$.
\end{lemma}
\begin{proof}
Cancel the nonzero factor $\rho_0 g a^3$ from both sides, leaving
$a/\sqrt{2} = \Delta h/4$; hence $a = \Delta h\, \sqrt{2}/4 =
\Delta h/(2\sqrt{2})$ using $\sqrt{2}^{\,2} = 2$.
\end{proof}

\begin{lemma}[Numerical readout at $\Delta h = 1.41$]
\label{lem:IPhO2026Problems_problem_IPhO_2026_1_A_1:numerical_value}
\lean{IPhO2026_1_A_1.numerical_value}
\uses{lem:IPhO2026Problems_problem_IPhO_2026_1_A_1:side_length_eq_delta_h_over}
For the design value $\Delta h = 1.41$ and $a = \Delta h/(2\sqrt{2})$, the
final readout is $a = 0.50 \;\lor\; |a - 0.50| < 1/200$.
\end{lemma}
\begin{proof}
The closed form gives $1.41/(2\sqrt{2}) = 0.4984\ldots$; the rational bounds
$1.414 < \sqrt{2} < 1.415$ place this value within $1/200$ of $0.50$.
\end{proof}

\begin{lemma}[Balance-law consistency bridge]
\label{lem:IPhO2026Problems_problem_IPhO_2026_1_A_1:torque_balance_contract}
\lean{IPhO2026_1_A_1.torque_balance_contract}
\uses{def:IPhO2026Problems_problem_IPhO_2026_1_A_1:HydrostaticGateSetup, def:IPhO2026Problems_problem_IPhO_2026_1_A_1:IsCriticalTorqueBalance}
The torque-balance existential bundled in the gate setup, re-expressed with
the derived magnitudes: a critical torque balance holds between the pressure
couple magnitude and the net immersed weight.
\end{lemma}
\begin{proof}
Project the torque-balance field of the setup bundle and identify its force
and couple with the defined magnitudes; no new physical input enters.
\end{proof}

\begin{theorem}[T1-A1 target: the gate side length]
\label{thm:IPhO2026Problems_problem_IPhO_2026_1_A_1:hydrostatic_gate_side_length_a_target}
\lean{IPhO2026_1_A_1.hydrostatic_gate_side_length_a_target}
\uses{def:IPhO2026Problems_problem_IPhO_2026_1_A_1:HydrostaticGateSetup, def:IPhO2026Problems_problem_IPhO_2026_1_A_1:PhysicalParameters, def:IPhO2026Problems_problem_IPhO_2026_1_A_1:PressureMomentReadout, lem:IPhO2026Problems_problem_IPhO_2026_1_A_1:net_immersed_weight_eq, lem:IPhO2026Problems_problem_IPhO_2026_1_A_1:weight_lever_arm_eq, lem:IPhO2026Problems_problem_IPhO_2026_1_A_1:restoring_moment_eq, lem:IPhO2026Problems_problem_IPhO_2026_1_A_1:pressure_couple_eq, lem:IPhO2026Problems_problem_IPhO_2026_1_A_1:critical_balance_eq, lem:IPhO2026Problems_problem_IPhO_2026_1_A_1:side_length_eq_delta_h_over, lem:IPhO2026Problems_problem_IPhO_2026_1_A_1:numerical_value, lem:IPhO2026Problems_problem_IPhO_2026_1_A_1:torque_balance_contract}
For the hydrostatic gate of Figure 1a in its critical (maximum-$\Delta h$)
configuration, with $\Delta h = 1.41$ m and the moment balance about O
holding, the side length of the cubic block is
$a = \Delta h/(2\sqrt{2})$ with $|a - 0.50| < 1/200$ — the official answer
$0.50$ m at the stated precision, appearing here conclusion-side only.
\end{theorem}
\begin{proof}
Positivity of $\rho_0$, $a$, $\Delta h$, $g$ comes from the parameter regime
inside the setup bundle. The consistency bridge re-expresses the bundled
balance as restoring moment equals pressure couple; Steps 5 and 7 turn this
into the scalar balance equation and solve it to $a = \Delta h/(2\sqrt{2})$.
The numerical readout at $\Delta h = 1.41$ then yields the neighbourhood of
$0.50$; the equality alternative trivially implies the bound since
$0 < 1/200$.
\end{proof}
